\section{Large grids and dispersive materials on the M5 Max}\label{sec:large}

Section~\ref{sec:studio} measured the build of patches 0001--0043 (cg33, commit 13350ca) on the M5 Max Mac Studio. On
2026-11-17 the same machine was used for upstream's own benchmark, \code{benchmark/bench.py}, at three sizes, and the
fork's CPU engine was behind upstream's new \code{cpu} engine at the largest one. This section describes what was done
about it: patches 0044 to 0049 (six commits, written 2026-11-21 to 2026-11-24 and committed 2026-11-22 to 2026-11-24, d918386 to 8af6c4c, snapshots cg34a to cg35c;
the fork's head is now 8af6c4c, v0.37.0-67-g8af6c4c). The installed build on the author's Mac Studio is still cg33; the
head is not installed. All numbers are from the M5 Max (6 ``Super'' and 12 ``Performance'' cores, 36\,GB, macOS~27.0) and
form part of the set of Section~\ref{sec:studio}; none is compared with a number of the M5 or the M5 Pro. The records
are \code{results/upstream\_bench\_studio\_2026-11-17.txt}, \code{results/large\_grid\_cg34\_2026-11-22.txt} and
\code{results/dispersive\_cg35\_2026-11-23.txt}.

\paragraph{Conditions of every table.} Upstream's \code{bench.py} with an odd number of mesh lines (101, 201 and 311), so that
every case has a source: with the even counts of upstream's own meshes (100 and 200 lines) no mesh line passes through the
origin and the ideal, pml and dispersive cases run with all fields zero (such runs are not a test of correctness). The sizes are S 1.03\,M, M 8.12\,M and L 30.08\,M cells (the cylindrical case 29.51\,M at L);
six cases (ideal, pml, inhomogeneous, dispersive, antenna, cylindrical); 6, 12 and 18 threads; 3 runs of 3\,s a point
(2\,s in the last slot); the builds run one after the other inside each point, run for run, and the table shows the best
rate over the thread counts, the median of the runs. ``master'' is upstream's commit 1d710c2 (with CSXCAD fdff4b5), its
\code{cpu} engine tiling along x only, and ``upstream head'' is ccb4472 (CSXCAD 7a24c75), tiling along two axes. The fork's
rows are in the CPU engine (\code{multithreaded}). The other processes used a median of 0.1 to 0.2 cores in every slot
(5\,s samples, the GPU 0\,\% outside the GPU parts; the loads of each part are in the records). In every comparison of
the fork's builds the probes are bit-identical (cg34c against cg33 on 54 of 54 rows, cg35a, cg35b and cg35c against their
predecessors on every row); the upstream builds differ from the fork, which is identical to v0.37.0, by $10^{-7}$ to
$3\times10^{-5}$ relative, and upstream's head differs on the cylindrical case by $7\times10^{-4}$ to
$2\times10^{-2}$, beyond the tolerance of \code{bench.py} (which side is right was not examined). The fork's Metal engine has no column: every case of the tool uses the
steady-state detection, which that engine runs on its CPU loops.

\subsection{Where cg33 stood}
Table~\ref{tab:largebefore} is the table of 2026-11-17: cg33 against master's \code{cpu} engine, paired. The fork led
at 1\,M cells (2.5 to 2.8$\times$ on five cases) and at 8\,M cells (1.10 to 1.29$\times$), was behind at 30\,M on every case
(0.83 to 0.86$\times$, cylindrical 0.99$\times$), and was at 0.48$\times$ on the dispersive case there. Master's \code{cpu}
engine held its rate from 8 to 30\,M cells on the ideal case (4\,320 to 4\,504\,\MC) while the fork's fell (4\,755 to 3\,755).

\begin{table}[H]\centering\small
\begin{tabular}{@{}lrrrrr@{}}
\toprule
 & \multicolumn{1}{c}{1\,M} & \multicolumn{1}{c}{8\,M} & \multicolumn{3}{c}{30\,M} \\
\cmidrule(lr){2-2}\cmidrule(lr){3-3}\cmidrule(l){4-6}
case & fork / master & fork / master & fork & master & fork / master \\
\midrule
ideal & 2.64 & 1.10 & 3\,755 & 4\,504 & 0.83 \\
pml & 2.65 & 1.13 & 2\,692 & 3\,143 & 0.86 \\
inhomogeneous & 2.67 & 1.13 & 2\,686 & 3\,142 & 0.85 \\
dispersive & 1.50 & 0.66 & 1\,378 & 2\,879 & 0.48 \\
antenna & 2.54 & 1.13 & 2\,511 & 2\,950 & 0.85 \\
cylindrical & 2.78 & 1.29 & 3\,860 & 3\,894 & 0.99 \\
\bottomrule
\end{tabular}
\caption{M5 Max, 2026-11-17: the fork's CPU engine (cg33, 13350ca) against upstream master's \code{cpu} engine (1d710c2),
best rate over 6, 12 and 18 threads, medians of 3 runs of 3\,s, interleaved; ratios are the fork's rate over master's; the
two columns of rates are in \MC. A source in every case; the other processes used a median of 0.10 to 0.16 cores.
\code{results/upstream\_bench\_studio\_2026-11-17.txt}.}\label{tab:largebefore}
\end{table}

\subsection{The cause}
The fused sweep (Section~\ref{sec:cg29}) walks each thread's x-slab in whole y-z planes and reuses the line behind. Once
four planes of voltages and currents no longer fit a thread's share of the L2 cache, that line has left the cache. On the
plain box at 18 threads the fork's rate falls from 7\,477\,\MC\ at 2\,M cells to 4\,803 at 8\,M, 4\,080 at 16\,M and 3\,778 at 30\,M;
master's, from 3\,206 to 4\,548 over the same range, crosses it at 12\,M cells. There is no sharp knee: a slope that
starts where four planes pass about 1\,MB. A timed test of master's own switch (blocked against flat, 30\,M cells) showed what its lead
consists of: on the plain box the temporal blocking is worth 1.22 to 1.35$\times$, on the pml case 1.02 to 1.08$\times$, so master's lead on the plain box
is its tiles, and on the pml case it is its flat kernel. The fork's opt-in tiled sweep (patch 0040, Section~\ref{sec:tiled}) already ran the
plain box at 30\,M cells at 5\,764\,\MC\ (against master's 4\,356 and the fused sweep's 3\,731 in the same slot), but
with a PML it ran at 0.26 to 0.79 of the fused sweep, because its hook stage gave whole x-lines to a few threads.

\subsection{Box hooks and the y-banded fused sweep (patches 0044 to 0046)}
\paragraph{Box hooks (0044, d918386).} Extensions get hooks that take an x range and a y range (\code{BOX\_HOOKS}; six
hooks, for the voltage and current updates, Pre and Post, and the two Apply calls). When every extension with a hook
declares them, a task of the tiled sweep runs Pre, the update, Post and Apply on each thread's own y lines behind one
barrier. Patch 0044 gives them to the UPML, the excitation, the lumped RLC elements and the steady-state detection;
\code{OPENEMS\_TBLOCK\_BOX=0} keeps the old hook stages. With them the tiled sweep with a PML at 30\,M cells and 18
threads ran at 1.15 to 1.28$\times$ the fused sweep (the hook stages: 0.32$\times$). It stays opt-in: on the plain box
at 18 threads it is unsteady, and with a
PML it lost at 8\,M cells.

\paragraph{The y-banded fused sweep (0045, e05d419).} The box hooks are exactly what is needed to cut the fused sweep along
y. Every block of a thread's x-slab, except the slab's last block, is cut into bands of $H$ y lines. Per band the
voltages of the block are updated (with the extensions' voltage box hooks), then the currents of the lines behind, on the
rows one line back in y (the first band takes row 0, the last band runs to the end), so the currents are skewed one row in y
as the fused sweep skews them one line in x. The planes that the sweep reuses stay in the thread's share of the L2 cache. The slab's last block
and the line before the slab keep whole planes, so the neighbour-flag protocol between threads is as before (Section~\ref{sec:cg29}).
The sweep is selected by \code{OPENEMS\_YBAND}: \code{0} (off), a number of lines (forced) or \code{auto}, the default
when the variable is not set. The automatic rule takes a thread's share of the cache as the smallest
\code{hw.perflevelN.l2cachesize} over \code{cpusperl2} of the performance levels the threads occupy (the number of MiB can
be set with \code{OPENEMS\_YBAND\_CACHE}); it uses bands when four y-z planes exceed that share, with $H$ equal to half
the share over four lines (at least 4), and it is off when two bands would not fit the grid. All or nothing: an
extension with hooks but without box hooks keeps whole planes.
On the M5 Max the share is 2.67\,MiB up to 6 threads (the Super cluster's 16\,MiB over 6 cores) and 1.33\,MiB above (the
8\,MiB clusters), so $311^3$ cells get bands of 93 lines at 6 threads and of 46 lines at 12 and 18, $201^3$ cells get
none at 6 threads and 71 lines at 12 and 18, and $101^3$ cells get none. At 30\,M cells and 18 threads the band height
matters little near the automatic value (Table~\ref{tab:largeband}).

\begin{table}[H]\centering\small
\begin{tabular}{@{}lrrrrr@{}}
\toprule
 & fused & 16 lines & 32 lines & auto (46) & 96 lines \\
\midrule
ideal & 3\,738 & 5\,171 & 5\,633 & 5\,661 & 4\,926 \\
pml & 2\,690 & 3\,766 & 4\,114 & 4\,136 & 3\,765 \\
\bottomrule
\end{tabular}
\caption{M5 Max, 30\,M cells, 18 threads: the fork with y-bands of different heights (cg34b), in \MC, medians of the runs of
one slot of 2026-11-22. \code{results/large\_grid\_cg34\_2026-11-22.txt}, steps before the final run.}\label{tab:largeband}
\end{table}

\paragraph{Box hooks for the remaining extensions (0046, 8d513ba).} With patches 0044 and 0045 the bands covered only the
ideal and the pml cases. Patch 0046 adds box hooks for the Mur boundary (the face cells are sorted by y per x-line; a y-face
couples two rows, so no band and no thread's y share is cut between them), for the conducting sheets and the Lorentz, Drude and
Debye materials (a packed group that a y range cuts runs on a thread-local copy and writes back only its inside lanes) and
for the closed cylinder (the extension declares a wrap in y: the last band runs row 0 first and the block of x-line 0
runs in whole planes). With these, every extension of the six cases and of the two boards has box hooks. Boards are below the
automatic rule's threshold and did not change.

\subsection{Result: 30\,M cells}
Table~\ref{tab:largefinal} is the final benchmark of 2026-11-22 (a quiet slot of 61 minutes: 252 points, 756 runs, none failed).
At 30\,M cells the fork leads upstream's \code{cpu} engines on five cases of six: 1.32 to 1.48$\times$ master and 1.11
to 1.22$\times$ upstream's head, and 1.49 to 1.55$\times$ its former self. At 8\,M cells it is 1.09 to 1.28$\times$
cg33, 1.36 to 1.58$\times$ master and 1.17 to 1.26$\times$ the head on five cases, where cg33 was only level with the head (0.97 to 1.02$\times$). At 1\,M cells the bands are off and cg34 equals cg33 within 1\,\% on
all six cases, still 1.53 to 2.89$\times$ master and 1.26 to 2.19$\times$ the head. The gain is at 12 and 18 threads
(at 30\,M cells 1.45 to 1.61$\times$ cg33). At 6 threads, the Super cores that the bus does not hold back, the bands
change little (0.96 to 1.03$\times$ cg33 at 30\,M cells) and master's \code{cpu} engine leads the fork on the cases with
extension work: 0.94 to 0.96$\times$ on pml, inhomogeneous and antenna at 30\,M cells and 0.85 to 0.89$\times$ at 8\,M,
the dispersive case 0.71 to 0.73$\times$. The dispersive case is the one case of six that stays behind in this table; the
band lifts it 1.54$\times$ at 30\,M cells but it remains at 0.75 of master and 0.64 of the head.

\begin{table}[H]\centering\small
\begin{tabular}{@{}llrrrrrrr@{}}
\toprule
 & & \multicolumn{4}{c}{rate, \MC} & \multicolumn{3}{c}{cg34 over} \\
\cmidrule(lr){3-6}\cmidrule(l){7-9}
size & case & cg33 & cg34 & master & head & master & head & cg33 \\
\midrule
L, 30\,M & ideal & 3\,755 & 5\,719 & 4\,341 & 5\,096 & 1.32 & 1.12 & 1.52 \\
 & pml & 2\,702 & 4\,149 & 2\,996 & 3\,411 & 1.38 & 1.22 & 1.54 \\
 & inhomogeneous & 2\,691 & 4\,173 & 2\,976 & 3\,464 & 1.40 & 1.20 & 1.55 \\
 & dispersive & 1\,375 & 2\,121 & 2\,845 & 3\,334 & 0.75 & 0.64 & 1.54 \\
 & antenna & 2\,537 & 3\,936 & 2\,969 & 3\,437 & 1.33 & 1.15 & 1.55 \\
 & cylindrical & 3\,901 & 5\,797 & 3\,909 & 5\,231 & 1.48 & 1.11 & 1.49 \\
\midrule
M, 8\,M & ideal & 4\,758 & 6\,092 & 4\,179 & 4\,914 & 1.46 & 1.24 & 1.28 \\
 & pml & 2\,991 & 3\,669 & 2\,564 & 2\,921 & 1.43 & 1.26 & 1.23 \\
 & inhomogeneous & 3\,006 & 3\,684 & 2\,597 & 2\,965 & 1.42 & 1.24 & 1.23 \\
 & dispersive & 1\,903 & 2\,068 & 2\,764 & 3\,155 & 0.75 & 0.66 & 1.09 \\
 & antenna & 2\,916 & 3\,415 & 2\,505 & 2\,921 & 1.36 & 1.17 & 1.17 \\
 & cylindrical & 5\,384 & 6\,336 & 3\,999 & 5\,304 & 1.58 & 1.19 & 1.18 \\
\midrule
S, 1\,M & ideal & 7\,635 & 7\,632 & 2\,963 & 3\,477 & 2.58 & 2.19 & 1.00 \\
 & pml & 3\,755 & 3\,726 & 1\,392 & 1\,890 & 2.68 & 1.97 & 0.99 \\
 & inhomogeneous & 3\,876 & 3\,866 & 1\,400 & 1\,917 & 2.76 & 2.02 & 1.00 \\
 & dispersive & 2\,796 & 2\,808 & 1\,836 & 2\,235 & 1.53 & 1.26 & 1.00 \\
 & antenna & 3\,708 & 3\,688 & 1\,420 & 1\,905 & 2.60 & 1.94 & 0.99 \\
 & cylindrical & 8\,427 & 8\,493 & 2\,935 & 4\,158 & 2.89 & 2.04 & 1.01 \\
\bottomrule
\end{tabular}
\caption{M5 Max, 2026-11-22: best rate over 6, 12 and 18 threads (18 on every point but cg33's dispersive L, at 6, and its
dispersive and antenna M, at 12) of the fork before (cg33, 13350ca) and with patches 0044 to 0046 (cg34, snapshot cg34c = 8d513ba), upstream master
(1d710c2) and upstream's head (ccb4472), CPU engines, in \MC, medians of 3 runs of 3\,s, interleaved; the last three columns are
the ratios of cg34's rate to the others. The other processes used a median of 0.20 cores (90th percentile 0.39), the
GPU 0\,\%. cg34 against cg33: probes identical on 54 of 54 rows, bands on in 30 of them. Cylindrical L is 29.5\,M cells.
\code{results/large\_grid\_cg34\_2026-11-22.txt}.}\label{tab:largefinal}
\end{table}

The patches were checked before they were timed: the six-model check against stock (18 of 18 rows identical), the Debye
check and the validation suite (15 of 15), raw fields and probes on 17 kinds of model for band heights of 1 to $N_y-1$
lines on 1 to 6 threads (1\,388 of 1\,388 identical), the closed cylinder (116 of 116) and the steady-state stop step with
Mur and with PML. The two-axis tiles that the plan allowed for were not needed: the bands reach master on the cases they
were aimed at.

\subsection{The dispersive pass (patches 0047 to 0049)}
About half of the dispersive case's cells are dispersive (a Lorentz block and a Debye block). After the bands it ran at 2\,121\,\MC\ at 30\,M cells with 6.4\,GiB of memory (6\,580\,MiB) and a setup of 13\,s, against
upstream's 2.0 to 2.2\,s. The packed hooks of patches
0042 and 0043 grouped four consecutive positions of an x-line, in a filled block four z of one column, that is four f4vectors
with one lane each: scalar gathers through twelve indices a group, and state and coefficients in one block that was
read and written back whole (about 132 bytes a Lorentz cell-step and 156 a Debye one, against 49 for a plain cell).
Three steps, each kept only if it passed a rule fixed before the run (Table~\ref{tab:largedisp}):

\begin{itemize}
\item \textbf{Vector groups (0047, d3374c9, cg35a).} A group is now the positions of one f4vector (a base offset and a lane
  mask; masked lanes keep their field bits) wherever a hook has at least 2.5 positions a group on average; thin structures
  such as the conducting sheets keep the position groups. The coefficients move to a table of distinct sets with a 16-bit
  or 32-bit index a group, apart from the state that is written. The per-position arrays are freed after packing. The
  arithmetic of every lane is unchanged. \code{OPENEMS\_ADE\_VEC=0} keeps the position groups (2 forces vector groups),
  \code{OPENEMS\_ADE\_COEF} picks the coefficient layout (both are test switches). Rule: 1.30$\times$ the dispersive case at 30\,M cells;
  result 1.41$\times$ (nothing else lost more than 2\,\%: antenna 0.99 to 1.02$\times$, the two boards 1.005$\times$).
\item \textbf{Debye transients in per-thread scratch (0048, f7cae41, cg35b).} A packed Debye group stored, beside the
  state of its poles, the pre-sum that Pre hands to Post and the correction that Post hands to Apply: 24 of $12(P+2)$ floats a group, written to the large
  block and read back every time step only to cross two calls. The engine now tells an extension after every plan
  (\code{RANGE\_LOCAL}, \code{SetRangeLocal}) whether Pre, Post and Apply of one range (an x-line block or an x-y box) run
  back to back on one thread, and the widest range. That holds for the fused sweep in line blocks, its y-bands with an
  open y direction and the box mode of the tiled sweep, not for the unfused plan, the tiled sweep's hook stages or bands of a
  closed y direction. While it holds, the transients live in a scratch slice of the calling thread and the blocks hold the pole
  state only (144 to 48 bytes written a group and step at one pole); a check of thread, width and call order stops the run if the engine breaks the grant. The
  Lorentz and Drude state is read again in the next step and is left as it was. \code{OPENEMS\_ADE\_SCRATCH=0} keeps the stored layout. Rule: 1.08$\times$; result 1.10$\times$ at 30\,M cells, with one loss, the dispersive case at 8\,M cells and 6 threads (no bands there) at
  0.97$\times$ (1\,920 to 1\,860\,\MC). The build with the stored layout runs 1\,864 there, so that cost is in the changed
  hook code, not in the scratch.
\item \textbf{Setup threads and setup memory (0049, 8af6c4c, cg35c).} The two operator extensions walked the whole grid on
  one thread and asked the material lookup for a \code{dynamic\_cast} a cell. They now scan per x-line on the operator's
  setup threads, the per-thread blocks are merged in x order (the serial order, so every array is the same byte for byte at
  every pool size), columns without a wanted primitive are skipped, and the material lookup remembers its cast per thread and setup
  generation. The setup arrays come from \code{mmap} on macOS (\code{tools/setup\_memory.h}), so freed setup memory leaves the
  process. \code{OPENEMS\_SETUP\_THREADS} sets the pool size and \code{OPENEMS\_SETUP\_TIMING=1} prints the phases. No
  per-timestep function changes. Rule: setup $-25$\,\% or memory $-10$\,\% and no rate more than 2\,\% lower; result $-63$\,\%
  and $-36$\,\%, rates 0.98 to 1.01 of before on ten points (the lowest, 6 threads at 30\,M cells, 1\,815 to 1\,782, inside
  the runs' ranges).
\end{itemize}

\begin{table}[H]\centering\small
\begin{tabular}{@{}lrrrrr@{}}
\toprule
 & \multicolumn{3}{c}{dispersive rate, \MC} & \multicolumn{2}{c}{30\,M cells} \\
\cmidrule(lr){2-4}\cmidrule(l){5-6}
step & 30\,M & 8\,M & 1\,M & setup, s & memory, MiB \\
\midrule
0047, vector groups & 2\,108 $\to$ 2\,962 & 2\,068 $\to$ 2\,976 & 2\,777 $\to$ 3\,942 & 13.0 $\to$ 13.2 & 6\,580 $\to$ 5\,857 \\
0048, Debye scratch & 2\,969 $\to$ 3\,266 & 2\,964 $\to$ 3\,156 & 3\,923 $\to$ 4\,054 & 13.9 $\to$ 13.2 & 5\,857 $\to$ 5\,857 \\
0049, setup & 3\,218 $\to$ 3\,233 & 3\,148 $\to$ 3\,172 & 4\,074 $\to$ 4\,074 & 13.3 $\to$ 4.9 & 5\,857 $\to$ 3\,737 \\
\bottomrule
\end{tabular}
\caption{M5 Max, 2026-11-23 and 2026-11-24: the dispersive case at 18 threads, before $\to$ after each patch, in \MC, one
paired slot a step (cg34 $\to$ cg35a, cg35a $\to$ cg35b, cg35b $\to$ cg35c), medians of 3 runs, upstream's builds
interleaved; setup time and resident memory at 30\,M cells; memory at 8\,M cells after 0049: 1\,742 to 1\,069\,MiB,
setup 3.63 to 1.55\,s. The other processes used a median of 0.05 to 0.2 cores. Antenna case
0.99 to 1.02$\times$ in all three steps (apart from one loaded 6-thread point at 1\,M cells, 1.08) and the two boards 0.996 to 1.014$\times$ for 0047 and 0048 (not timed for 0049). All fork rows identical to their
predecessor's. \code{results/dispersive\_cg35\_2026-11-23.txt}, three sections.}\label{tab:largedisp}
\end{table}

Table~\ref{tab:largedispfinal} is the comparison with upstream in the slot of the last step. The fork now leads master
on the dispersive case at all three sizes and is level with upstream's head at 8\,M cells (1.04$\times$) and just behind
at 30\,M (0.97$\times$). In the slot of the second step master ran 2\,883 at 30\,M cells and 2\,803 at 8\,M, which makes the ratio
1.13 at both sizes (1.13 to 1.20 at 30\,M and 1.13 to 1.39 at 8\,M over the two slots: master's runs spread widely in
the last one, 2\,073 to 2\,852 at 8\,M cells), and in that slot the head ran 1.00$\times$ at 8\,M and 0.97$\times$ at 30\,M.

\begin{table}[H]\centering\small
\begin{tabular}{@{}lrrrrrrrr@{}}
\toprule
 & \multicolumn{3}{c}{rate, \MC} & \multicolumn{2}{c}{fork over} & \multicolumn{3}{c}{setup, s / memory, MiB} \\
\cmidrule(lr){2-4}\cmidrule(lr){5-6}\cmidrule(l){7-9}
size & fork & master & head & master & head & fork & master & head \\
\midrule
30\,M & 3\,233 & 2\,700 & 3\,343 & 1.20 & 0.97 & 4.90 / 3\,737 & 2.08 / 3\,961 & 2.26 / 2\,680 \\
8\,M & 3\,172 & 2\,284 & 3\,048 & 1.39 & 1.04 & 1.55 / 1\,069 & 0.58 / 1\,231 & 0.64 / 921 \\
1\,M & 4\,074 & 1\,895 & 2\,248 & 2.15 & 1.81 & 0.25 / 164 & 0.10 / 158 & 0.10 / 117 \\
\bottomrule
\end{tabular}
\caption{M5 Max, 2026-11-24 (08:39 to 08:45): the dispersive case, the fork with patches 0001 to 0049 (cg35c, 8af6c4c)
against upstream master (1d710c2) and upstream's head (ccb4472), CPU engines, 18 threads, medians of 3 runs of 2\,s,
interleaved; rates in \MC; setup time and resident memory of the same runs. The other processes used a median of 0.16 cores. Master's runs spread widely in this slot at 30\,M cells (2\,647 to
2\,884) and at 8\,M cells (2\,073 to 2\,852). Fork probes identical to cg35b's; the upstream builds differ by
$3\times10^{-6}$ to $3\times10^{-5}$. \code{results/dispersive\_cg35\_2026-11-23.txt}, third section.}\label{tab:largedispfinal}
\end{table}

The checks before each timing were the same kind as for the bands: 630, 501 and 470 box configurations for the three
patches, the dispersive check against stock at 50\,k and 1\,M cells and on the closed cylinder, the six-model check
(18 of 18 identical), the Debye check, the validation suite (15 of 15) and planted faults (for the last patch six of eight were caught; the two others need an index pattern no test model has). Patch 0049 changes no per-timestep function.

\subsection{What stays open}
\begin{itemize}
\item The dispersive case at 30\,M cells is at 0.97 of upstream's head; upstream sets it up in 2.1 to 2.3\,s against the fork's 4.9 and its head
  uses 2\,680\,MiB against 3\,737.
\item At 6 threads master (1d710c2) still leads the fork on the cases with extension work at 8 and 30\,M cells: 0.85 to
  0.96$\times$, the dispersive case 0.90 to 0.96$\times$ after the pass. The bands do not help there, and the Super cores are not
  held back by the bus.
\item The two boards are below the automatic rule's threshold (their grids are small); the bands are not a gain for them, and
  they were measured unchanged (0.996 to 1.014$\times$ for patches 0047 and 0048).
\item The tiled sweep (patch 0040) stays opt-in; with box hooks it gains on pml at 30\,M cells but is unsteady at 18 threads and loses at 8\,M. A choice at run time between the sweeps is not built.
\item No Metal column in these tables: the steady-state detection that every case of \code{bench.py} uses has no GPU version; with
  it, the fork's Metal engine runs the CPU loops of that extension and is not measured here.
\item Upstream's head differs from the fork and from master on the cylindrical case beyond \code{bench.py}'s tolerance;
  which result is right was not examined.
\end{itemize}
